% !TeX root = ../main.tex
\chapter{Related Work}
\label{ch:related-work}
\thesisplaceholder{Synthesize verified primary literature and explain the unresolved research gap.}

\section{Clinical and Biological Background}
% 整理任務相關的生理、病理或照護流程；區分普遍知識與特定研究情境。
% 指出標籤或量測與真實臨床／生物狀態的差距。

\section{Data and Evaluation in Prior Studies}
% 依族群、資料來源、切分、終點、外部驗證與限制比較研究。
% 不能直接比較任務、族群或評估設定不相容的 headline scores。

\section{Methods and Baselines}
% 先整理現行實務、簡單統計方法與強基準，再介紹相關 ML 方法。
% 解釋選擇理由與方法假設；勿以新穎架構本身作為充分研究貢獻。

\section{Evidence Gap and Study Positioning}
% 逐項核對文獻是否支持研究缺口；區分未研究、證據不足與已得到負結果。
% 每個關鍵主張都應有實際閱讀並核實的來源。
